
\begin{table*}[htbp!]
\centering
%\small
%    \setlength{\tabcolsep}{2pt}
\begin{threeparttable}[c]
\caption{Characteristic parameters of the DCN (3-2) hyperfine fits for the non-background subtracted}
%(using global measurements)}
\label{tabappendix:dcnfits_average}
 \begin{tabular}{lllllllllllll}
\hline \noalign {\smallskip}
%Protocluster  & \multicolumn{3}{c}{spw0 }  & \multicolumn{3}{c}{spw1}  & \multicolumn{3}{c}{spw2 }
%    & \multicolumn{3}{c}{spw3} & \\ %& \lbol & \tdust 
%cloud name\tnote{1}    & min & max & rest   & min & max & rest  & min & max & rest  & min & max & rest   \\
Protocluster         & Number\tnote{1}  & Number\tnote{2}   & Detection\tnote{2} & Velocity\tnote{3} & \vlsr\tnote{3}   & \vlsr\tnote{3}    & Linewidth\tnote{3}  & Linewidth\tnote{3}   &  Linewidth\tnote{3}  \\ 
cloud name & of Cores & detected    & rate [\%]   & range    & mean   & std     & mean  & median &  std \\
&  &(complex)        &  (complex)      & [km~s$^{-1}$]   & [km~s$^{-1}$] & [km~s$^{-1}$]  & [km~s$^{-1}$]& [km~s$^{-1}$] & [km~s$^{-1}$] \\

\hline \noalign {\smallskip}
\multicolumn{10}{c}{Young} \\
G327.29 & 32 & 18 (23)&  56 (72)\% & 5.3 & -44.9 & 1.9  & 1.7 & 1.7 & 0.5 \\
G328.25 & 11 & 3 (4)&  27 (36)\% & 1.4 & -43.1 & 0.6  & 1.9 & 1.6 & 0.9 \\
G338.93 & 42 & 21 (29)&  50 (69)\% & 6.9 & -60.9 & 2.0  & 2.1 & 2.0 & 0.7 \\
G337.92 & 22 & 10 (17)&  45 (77)\% & 4.0 & -39.2 & 1.3  & 2.1 & 2.0 & 1.1 \\
W43-MM1 & 70 & 32 (51)&  46 (73)\% & 7.5 & 97.2 & 2.1  & 1.9 & 1.6 & 1.1 \\
W43-MM2 & 40 & 23 (30)&  58 (75)\% & 4.3 & 90.8 & 1.3  & 1.9 & 1.6 & 0.9 \\
Average &  36 & 18(26) & 49 (71)\% &  4.9 &  & 1.8  & 1.9 & 1.7 & 0.9 \\
Total &  217 & 107 (154)  \\
\hline
\multicolumn{10}{c}{Intermediate} \\
G351.77 & 18 & 5 (14)&  28 (78)\% & 5.9 & -4.5 & 2.1  & 2.0 & 1.6 & 0.6 \\
G353.41 & 45 & 17 (22)&  38 (49)\% & 7.4 & -17.3 & 1.8  & 1.7 & 1.5 & 0.8 \\
G008.67 & 19 & 8 (13)&  42 (68)\% & 6.6 & 35.8 & 2.0  & 2.0 & 1.9 & 0.6 \\
W43-MM3 & 36 & 12 (18)&  33 (50)\% & 2.6 & 93.1 & 0.8  & 1.8 & 2.0 & 0.7 \\
W51-E & 23 & 12 (21)&  52 (91)\% & 11.1 & 56.8 & 3.9  & 2.0 & 1.6 & 1.0 \\
Average &  28 & 11(18) & 38 (62)\% &  6.7 &  & 2.8  & 1.8 & 1.6 & 0.8 \\
Total &  141 & 54 (88)  \\
\hline
\multicolumn{10}{c}{Evolved} \\
G333.60 &  52 & 36 (42) &  69 (81)\% & 10.0 & -47.8 & 1.9  & 1.7 & 1.7 & 0.5 \\
G010.62 &  42 & 28 (37) &  67 (88)\% & 9.3 & -2.5 & 2.1  & 1.4 & 1.3 & 0.6 \\
G012.80 &  46 & 38 (40) &  83 (87)\% & 6.4 & 36.1 & 1.4  & 1.4 & 1.2 & 0.6 \\
W51-IRS2 &  97 & 69 (84) &  71 (87)\% & 18.3 & 61.9 & 2.9  & 1.5 & 1.3 & 0.8 \\
Average &  59 & 43 (51) & 72 (86)\% &  11.0 &  & 2.5  & 1.5 & 1.3 & 0.7 \\
Total &  237 & 171 (203)  \\
\hline \noalign {\smallskip}
\end{tabular}
\begin{tablenotes}
\item[1] The number of continuum cores are taken from the smoothed core catalog from \citet{louvet22}. %We note that 2 cores overlap in W51-E and W51-IRS2 regions and 2 cores overlap in W43-MM2/3 regions (we note these in \cref{appendix:dcn_fitting}), however, as they are either non-detections or Complex-type spectra we do not explicitly assign them to a given field. 
\item[2] The total count/percentages in the brackets include both Single- and Complex-type DCN (3-2) detections. Where spectra that are described as Complex-type spectra and are not well fit by a single component. We provide the percentages and number of detections, but cores with a Complex-type spectrum are not included in the estimates of the Velocity range, \vlsr, or linewidth.
\item[3] Only Single-type DCN spectra (fit with a single component) and listed as Single in the Tables of the DCN fits (e.g. \cref{tab:coretables_g338}) are included in the estimates of the Velocity ranges, \vlsr and linewidth for each protocluster. 
%
%
\end{tablenotes}
\end{threeparttable}
\end{table*}